\chapter{تكامل لوبيغ}\label{ch:b3:lebesgue}

يشرّح تكامل ريمان \emph{مجموعة التعريف} إلى فترات صغيرة؛ أما
تكامل لوبيغ فيشرّح \emph{مجموعة القيم}: فلمكاملة $f$ نقيس
المجموعات $\{f > t\}$. ويبدو هذا التغيير بريئًا وهو ثوري.
فالنهايات والتكاملات، المتخاصمة أبدًا في نظرية ريمان (إذ يُطلب
التقارب المنتظم!)، تتصالح بواسطة ثلاث مبرهنات في التقارب ---
التقارب الرتيب، وفاتو، والتقارب المهيمن --- وفرضياتها ضعيفة إلى
حدّ يكاد يُخجل. ويبني هذا الفصل التكامل على \omterm{def:b3:measure:measure}{فضاء قياس} كيفي
$(X, \mathcal A,
\mu)$، ويبرهن على المبرهنات الثلاث، ويحسم العلاقة المضبوطة مع
تكامل ريمان (فالدالة المحدودة \omterm{def:b3:lebesgue:l1}{قابلة للمكاملة} بمعنى ريمان إذا
وفقط إذا كانت متصلة \omterm{def:b3:lebesgue:l1}{في كل مكان تقريبًا})، ويُصنّع اشتقاق
\omterm{thm:b3:lebesgue:paramcont}{التكاملات المتعلقة بوسيط} --- وهي التقنية التي تستعملها مسألة
نهاية الأسبوع لحساب $\int_0^\infty\frac{\sin x}x\,\dd x$ و $\int_\R \eu^{-x^2}\dd x$.

\section{الدوال القابلة للقياس}

\begin{definition}\label{def:b3:lebesgue:measurable}
ليكن $(X, \mathcal A)$ و $(Y, \mathcal B)$ فضاءين قابلين \omterm{def:b3:measure:measure}{للقياس}. يكون $f \colon X \to Y$ \emph{قابلًا
للقياس}\index{دالة قابلة للقياس} إذا كان $f^{-1}(B) \in \mathcal A$ من أجل كل $B \in
\mathcal B$.
ومن أجل الدوال الحقيقية (أو ذات القيم في $[-\infty,+\infty]$)، يحمل $Y = \R$
جبر بوريل من النمط $\sigma$، ويكفي التحقق من $f^{-1}(\intoo t{+\infty}) = \{f > t\} \in
\mathcal A$ من أجل كل
$t \in \R$: إذ تشكّل المجموعات الجيدة $\{B :
f^{-1}(B) \in \mathcal A\}$ جبرًا من النمط
$\sigma$ (لأن السوابق تتبادل مع عمليات المجموعات) يحتوي على
الأشعة المولِّدة (\cref{def:b3:measure:borel}، \cref{met:b3:measure:goodsets}).
\end{definition}

\begin{proposition}\label{prop:b3:lebesgue:stability}
(a) تركيب التطبيقات القابلة \omterm{def:b3:measure:measure}{للقياس} قابل \omterm{def:b3:measure:measure}{للقياس}؛ والتطبيقات
المتصلة قابلة \omterm{def:b3:measure:measure}{للقياس} بمعنى بوريل.
(b) إذا كانت $f, g \colon X \to \R$ قابلة \omterm{def:b3:measure:measure}{للقياس}، فكذلك $f + g$ و $fg$
و $\max(f,g)$ و $\abs f$ و $\lambda f$.
(c) إذا كانت $(f_n)$ قابلة \omterm{def:b3:measure:measure}{للقياس} بقيم في $[-\infty, +\infty]$، فإن
$\sup_nf_n$ و $\inf_nf_n$ و $\limsup f_n$ و $\liminf f_n$ قابلة
\omterm{def:b3:measure:measure}{للقياس}؛ وإذا كان $f_n \to f$ نقطةً نقطة، فإن $f$ قابلة \omterm{def:b3:measure:measure}{للقياس}.
\end{proposition}

\begin{proof}
(a) $(g\circ f)^{-1}(B) = f^{-1}(g^{-1}(B))$؛ ويعطي الاتصال القابليةَ \omterm{def:b3:measure:measure}{للقياس} عبر \omterm{def:b3:topology:topology}{المجموعات
المفتوحة} المولِّدة (\cref{pb:b3:measure:1}، السؤال 10، في صيغته العامة).
(b) التطبيق $(f, g) \colon X \to \R^2$ قابل \omterm{def:b3:measure:measure}{للقياس} من أجل جبر بوريل من النمط
$\sigma$ للفضاء $\R^2$ --- بالتحقق على الصناديق المفتوحة، وهي
مولِّدة (لأن مفتوحات $\R^2$ اتحادات قابلة للعدّ لصناديق ناطقة):
$(f,g)^{-1}(U\times V) = f^{-1}(U)\cap g^{-1}(V)$ --- ويكون $+, \times, \max$ متصلين $\R^2 \to \R$: ثم نركّب.
(c) $\{\sup f_n > t\} = \bigcup_n\{f_n > t\}$؛ $\inf = -\sup(-)$؛ $\limsup = \inf_N\sup_{n \geq N}$؛ والنهاية نقطةً نقطة هي $\limsup$
لنفسها.
\end{proof}

\begin{definition}\label{def:b3:lebesgue:simple}
\emph{الدالة البسيطة}\index{دالة بسيطة} هي \omterm{def:b3:lebesgue:measurable}{دالة قابلة للقياس} ذات
عدد منتهٍ من القيم: $s = \sum_{i=1}^n
c_i\,\mathbf 1_{A_i}$، حيث $A_i \in \mathcal A$ منفصلة و $c_i \geq
0$ (من أجل
النظرية غير السالبة). وتكاملها هو
\[
\int s\,\dd\mu = \sum_i c_i\,\mu(A_i) \in [0, +\infty]
\]
(بالاصطلاح $0\cdot\infty = 0$)؛ ولا تتعلق القيمة \omterm{def:b3:representations:rep}{بالتمثيل} المختار (بتنقيح
تجزئتين).
\end{definition}

\begin{theorem}[التقريب بالدوال البسيطة]
\label{thm:b3:lebesgue:approximation}
كل \omterm{def:b3:lebesgue:measurable}{دالة قابلة للقياس} $f \colon X \to [0, +\infty]$ هي النهاية نقطةً نقطة لمتتالية
\emph{متزايدة} من الدوال البسيطة:
\[
s_n = \sum_{k=1}^{n2^n} \frac{k-1}{2^n}\,
\mathbf 1_{\{\frac{k-1}{2^n} \leq f < \frac k{2^n}\}}
+ n\,\mathbf 1_{\{f \geq n\}} \nearrow f .
\]
\end{theorem}

\begin{proof}
كل $s_n$ بسيطة (لأن المجموعات سوابق لمجموعات بوريلية). وأما
الرتابة: فالانتقال من $n$ إلى $n+1$ يشطر كل مستوى ثنائي إلى
اثنين ولا ينقص القيمة المُسنَدة أبدًا (فنقطة تحقق $\frac{k-1}{2^n} \leq f(x) < \frac k{2^n}$ تأخذ
إما $\frac{2k-2}{2^{n+1}}$ وإما $\frac{2k-1}{2^{n+1}}$، وكلتاهما $\geq
\frac{k-1}{2^n}$؛ ويرتفع السقف $n$
أيضًا). وأما التقارب: فإذا كان $f(x) < \infty$، فمن أجل
$n > f(x)$ لدينا $f(x) - s_n(x) \leq
2^{-n}$؛ وإذا كان $f(x) = \infty$، فإن $s_n(x) = n \to \infty$.
\end{proof}

\section{التكامل ومبرهنات التقارب}

\begin{definition}\label{def:b3:lebesgue:integral}
من أجل $f \geq 0$ قابلة \omterm{def:b3:measure:measure}{للقياس}:
\[
\int f \,\dd\mu = \sup\Bigl\{\int s\,\dd\mu : s \text{ بسيطة},
\ 0 \leq s \leq f\Bigr\} \in [0, +\infty].
\]
وهو رتيب في $f$ بحكم البناء، ويمدّد الحالة البسيطة (فمن أجل $f$
بسيطة، يُبلَغ الحد الأعلى عند $f$: بمقارنة تكاملات الدوال
البسيطة عبر تنقيحات مشتركة).
\end{definition}

\begin{theorem}[التقارب الرتيب، بيبو ليفي]
\label{thm:b3:lebesgue:mct}
إذا كان $0 \leq f_n \nearrow f$ نقطةً نقطة (بدوال قابلة \omterm{def:b3:measure:measure}{للقياس})، فإن
\[
\int f_n\,\dd\mu \nearrow \int f\,\dd\mu .
\]
\index{مبرهنة التقارب الرتيب}
\end{theorem}

\begin{proof}
الدالة $f$ قابلة \omterm{def:b3:measure:measure}{للقياس} (\cref{prop:b3:lebesgue:stability}(c)) ويتزايد $\int f_n$ إلى
$L \leq \int f$ ما (بالرتابة). وبالعكس، لنثبّت \omterm{def:b3:lebesgue:simple}{دالة بسيطة} $s = \sum c_i\mathbf 1_{A_i} \leq f$
وعددًا $\theta \in (0,1)$؛ فالمجموعات $E_n = \{f_n \geq \theta s\}$ قابلة \omterm{def:b3:measure:measure}{للقياس} وتتزايد إلى $X$
(فحيث $s(x) > 0$: $f(x)
\geq s(x) > \theta s(x)$، ومنه في نهاية المطاف $f_n(x) \geq \theta
s(x)$؛ وحيث
$s(x) = 0$: بداهةً). عندئذٍ
\[
\int f_n \geq \int_{E_n}\theta s\,\dd\mu
= \theta\sum_i c_i\,\mu(A_i \cap E_n)
\xrightarrow[n\to\infty]{} \theta\sum_ic_i\,\mu(A_i)
= \theta\int s
\]
بالاتصال من الأسفل (\cref{prop:b3:measure:basics}(c)). ومنه $L \geq \theta\int s$ من أجل كل
$\theta < 1$ ومن أجل كل \omterm{def:b3:lebesgue:simple}{دالة بسيطة} $s
\leq f$: أي $L \geq \int f$.
\end{proof}

\begin{corollary}\label{cor:b3:lebesgue:additivity}
من أجل $f, g \geq 0$ قابلتين \omterm{def:b3:measure:measure}{للقياس} ومن أجل $c \geq 0$: $\int(f + g) =
\int f + \int g$
و $\int cf = c\int f$؛ ومن أجل متسلسلة من الدوال غير السالبة القابلة \omterm{def:b3:measure:measure}{للقياس}،
$\int\sum_nf_n =
\sum_n\int f_n$.
\end{corollary}

\begin{proof}
من أجل الدوال البسيطة، تكون الجمعية حسابًا على تنقيح مشترك. وفي
العموم نأخذ $s_n \nearrow f$ و $t_n \nearrow g$ (\cref{thm:b3:lebesgue:approximation}):
فنجد $s_n + t_n \nearrow f +
g$، وتنقل مبرهنة التقارب الرتيب الجمعيةَ إلى النهاية.
وأما عبارة المتسلسلة فهي مبرهنة التقارب الرتيب مطبَّقةً على
المجاميع الجزئية.
\end{proof}

\begin{theorem}[مبرهنة فاتو المساعدة]\label{thm:b3:lebesgue:fatou}
من أجل دوال قابلة \omterm{def:b3:measure:measure}{للقياس} $f_n \geq 0$:
\[
\int \liminf_n f_n \,\dd\mu \;\leq\; \liminf_n \int
f_n\,\dd\mu .
\]
\index{مبرهنة فاتو المساعدة}
\end{theorem}

\begin{proof}
لتكن $g_N = \inf_{n\geq N}f_n$: فهي قابلة \omterm{def:b3:measure:measure}{للقياس}، و $0 \leq g_N \nearrow
\liminf f_n$، ويكون $g_N \leq f_n$
من أجل كل $n \geq N$، ومنه $\int g_N \leq \inf_{n \geq N}\int f_n$. ونطبّق مبرهنة التقارب الرتيب
على الطرف الأيسر: $\int\liminf f_n = \lim_N\int g_N \leq
\lim_N\inf_{n\geq N}\int f_n = \liminf\int f_n$.
\end{proof}

\begin{definition}\label{def:b3:lebesgue:l1}
تكون \omterm{def:b3:lebesgue:measurable}{الدالة القابلة للقياس} $f \colon X \to \R$ (أو ذات القيم في $\C$)
\emph{قابلة للمكاملة}\index{دالة قابلة للمكاملة} إذا كان $\int\abs
f\,\dd\mu < \infty$؛
وعندئذٍ $\int f = \int f^+ - \int f^-$ (بالجزأين الموجب والسالب؛ وبالجزأين الحقيقي
والتخيلي في الحالة العقدية). والتكامل خطي على الدوال القابلة
للمكاملة (بالتفكيك وإعادة التركيب على الأجزاء الموجبة؛ وتُختزل
الحالة العقدية إلى الحقيقية) ويحقق $\abs{\int f} \leq
\int\abs f$ (في الحالة الحقيقية:
$\pm\int f = \int(\pm f) \leq
\int\abs f$؛ وفي الحالة العقدية: بالضرب في ثابت أحادي المعيار لجعل
التكامل حقيقيًا). ونقول إن خاصية تتحقق \emph{في كل مكان تقريبًا}
\index{في كل مكان تقريبا} إذا لم تخفق إلا على مجموعة معدومة
\omterm{def:b3:measure:measure}{القياس} بالمعنى $\mu$؛ وتغيير $f$ على مجموعة معدومة \omterm{def:b3:measure:measure}{القياس} لا
يغيّر أي تكامل (لأن الفرق مهيمَن عليه بالمقدار $\infty\cdot\mathbf
1_N$، وتكامله $0$).
\end{definition}

\begin{theorem}[التقارب المهيمن]\label{thm:b3:lebesgue:dct}
ليكن $f_n \to f$ \omterm{def:b3:lebesgue:l1}{في كل مكان تقريبًا}، مع $\abs{f_n} \leq g$ \omterm{def:b3:lebesgue:l1}{في كل مكان تقريبًا}
من أجل دالة \emph{\omterm{def:b3:lebesgue:l1}{قابلة للمكاملة}} مثبَّتة $g$. عندئذٍ تكون $f$
\omterm{def:b3:lebesgue:l1}{قابلة للمكاملة} و
\[
\int f_n\,\dd\mu \longrightarrow \int f\,\dd\mu,
\qquad\text{بل}\quad \int\abs{f_n - f}\,\dd\mu \to 0 .
\]
\index{مبرهنة التقارب المهيمن}
\end{theorem}

\begin{proof}
نطرح مجموعة معدومة \omterm{def:b3:measure:measure}{القياس} لنجعل الفرضيات نقطةً نقطة. ولدينا
$\abs f
\leq g$: أي إن $f$ \omterm{def:b3:lebesgue:l1}{قابلة للمكاملة}. وتحقق الدوال $h_n = 2g - \abs{f_n
- f} \geq 0$ الشرطَ
$\liminf h_n = 2g$؛ وتعطي مبرهنة فاتو
\[
\int 2g \leq \liminf\int\bigl(2g - \abs{f_n - f}\bigr)
= \int 2g - \limsup\int\abs{f_n - f},
\]
ومنه $\limsup\int\abs{f_n - f} \leq 0$ (والطرح مشروع: $\int 2g < \infty$). وأخيرًا $\abs{\int f_n - \int f} \leq
\int\abs{f_n - f} \to 0$.
\end{proof}

\begin{method}\label{met:b3:lebesgue:convergence}
أمام $\lim_n\int f_n$: جرّب، بالترتيب --- (1) هل المتتالية
رتيبة (أو متسلسلة بحدود غير سالبة)؟ فمبرهنة التقارب الرتيب، دون
حاجة إلى القابلية للمكاملة. (2) هل يوجد مهيمِن واحد قابل
للمكاملة $g \geq \abs{f_n}$، يُوجَد بحدود فجّة (بأخذ «$\sup_n$» للتقديرات)؟
فمبرهنة التقارب المهيمن. (3) لا هيمنة ولا رتابة؟ فما تزال
مبرهنة فاتو تحدّ أحد الطرفين، وقد يخفق التساوي فعلًا: فالنتوء
الهارب $f_n = n\mathbf
1_{\intoo0{1/n}}$ يحقق $\int f_n = 1$ لكن $f_n \to 0$ \omterm{def:b3:lebesgue:l1}{في كل مكان
تقريبًا}. فالهيمنة هي بالضبط ما يمنع الكتلة من الهروب إلى
اللانهاية، رأسيًا أو أفقيًا.
\end{method}

\section{ريمان في مقابل لوبيغ}

\begin{theorem}[محك لوبيغ]
\label{thm:b3:lebesgue:riemann}
لتكن $f \colon \intcc ab \to \R$ \emph{محدودة}. عندئذٍ تكون $f$ \omterm{def:b3:lebesgue:l1}{قابلة للمكاملة} بمعنى
ريمان إذا وفقط إذا كانت $f$ متصلة \omterm{def:b3:lebesgue:l1}{في كل مكان تقريبًا} بالمعنى
$\lambda$؛ وفي تلك الحالة تكون $f$ \omterm{def:b3:lebesgue:l1}{قابلة للمكاملة} بمعنى لوبيغ
ويتطابق التكاملان.\index{تكامل ريمان في مقابل تكامل لوبيغ}
\end{theorem}

\begin{proof}
من أجل تجزئة $\sigma = (a = x_0 < \dots < x_N = b)$، ليكن $L_\sigma$ و $U_\sigma$ الدالتين
الدرجيتين المساويتين، على كل $\intoo{x_{i-1}}{x_i}$، للمقدارين $m_i = \inf_{[x_{i-1}, x_i]}f$
و $M_i = \sup$؛ ومجاميع داربو هي تكاملاهما (فريمان ولوبيغ
يتوافقان على الدوال الدرجية، وكلاهما يعطي $\sum
m_i\Delta x_i$). ونأخذ متتالية
من التجزئات $\sigma_n$، كلٌّ منها تنقيح لسابقتها، خطوتها
$\to 0$، بمجاميع داربو تتقارب إلى تكاملَي داربو الأدنى والأعلى
للدالة $f$. وتجعل التنقيحاتُ $L_{\sigma_n}$ غير متناقصة
و $U_{\sigma_n}$ غير متزايدة نقطةً نقطة خارج المجموعة القابلة للعدّ
$D$ لجميع نقط التقسيم؛ ولنسمِّ النهايتين $\ell$ و $u$ (وهما
قابلتان \omterm{def:b3:measure:measure}{للقياس}، \cref{prop:b3:lebesgue:stability}). ومن أجل $x \notin
D$، وبكتابة $I_n(x)$
للفترة المفتوحة من $\sigma_n$ الحاوية على $x$: $\ell(x) = \sup_n\inf_{I_n(x)}f$ و $u(x) =
\inf_n\sup_{I_n(x)}f$؛
وبما أن الخطوات تتقلّص إلى $0$، فهاتان هما \emph{الغلافان}
الأدنى والأعلى للدالة $f$ عند $x$ --- ويكون $u(x) - \ell(x)$
تذبذبَ $f$ عند $x$ --- بحيث يكون \emph{$\ell(x) = u(x)$ إذا
وفقط إذا كانت $f$ متصلة عند $x$}. وبمبرهنتَي التقارب
الرتيب/المهيمن (لأن الدالة محدودة والفترة منتهية):
\[
\int_{\intcc ab}\ell\,\dd\lambda = \lim_n\int L_{\sigma_n}
= \underline{\int}f,
\qquad
\int_{\intcc ab}u\,\dd\lambda = \overline{\int}f .
\]
فتكون $f$ \omterm{def:b3:lebesgue:l1}{قابلة للمكاملة} بمعنى ريمان $\iff$ $\underline\int f =
\overline\int f$ $\iff$
$\int(u - \ell) = 0$ $\iff$ $u = \ell$ \omterm{def:b3:lebesgue:l1}{في كل مكان تقريبًا} ($u - \ell \geq 0$؛
\cref{exo:b3:lebesgue:5}) $\iff$ تكون $f$ متصلة \omterm{def:b3:lebesgue:l1}{في كل مكان تقريبًا}. وفي تلك
الحالة $\ell \leq f \leq u$ مع $\ell =
u$ \omterm{def:b3:lebesgue:l1}{في كل مكان تقريبًا}: أي إن $f$ تساوي
$\ell$ القابلة \omterm{def:b3:measure:measure}{للقياس} \omterm{def:b3:lebesgue:l1}{في كل مكان تقريبًا}، ومنه فهي قابلة \omterm{def:b3:measure:measure}{للقياس}
بمعنى لوبيغ (بتمام $\lambda$) مع $\int f\,\dd\lambda = \int\ell\,\dd\lambda = \underline\int f =
\int_a^bf$.
\end{proof}

\begin{example}\label{ex:b3:lebesgue:riemannexamples}
الدالة $\mathbf 1_\Q$ غير متصلة في أي نقطة: فهي غير \omterm{def:b3:lebesgue:l1}{قابلة
للمكاملة} بمعنى ريمان --- لكنها تافهة بمعنى لوبيغ: $\int\mathbf 1_\Q\,\dd\lambda =
\lambda(\Q) = 0$. ودالة
توماي (\,$\frac1q$ عند الأعداد الناطقة $\frac pq$، و $0$ فيما
عدا ذلك) متصلة عند الأعداد الصمّاء بالضبط: فهي \omterm{def:b3:lebesgue:l1}{قابلة للمكاملة}
بمعنى ريمان وتكاملها $0$. وأما التكاملات \emph{المعتلّة} بمعنى
ريمان فمفهوم مختلف: إذ يتقارب $\int_0^\infty\frac{\sin x}x\,\dd x$ بوصفه نهايةً للتكاملات
$\int_0^A$ (وتحسبه مسألة نهاية الأسبوع $= \frac\pi2$)، لكن
$\frac{\sin x}x \notin L^1(\intoo0{+\infty})$: فالتكامل المطلق يتباعد كالمتسلسلة التوافقية
(\cref{exo:b3:lebesgue:6}). فنظرية لوبيغ تقايض التقارب الشرطي بمبرهنات
نهايات متينة.
\end{example}

\section{التكاملات المتعلقة بوسيط}

في كل ما يلي، $(X, \mathcal A, \mu)$ \omterm{def:b3:measure:measure}{فضاء قياس}، و $T$ فضاء متري (وهو الوسيط)،
و $f \colon T \times X \to \C$ مع كون $f(t, \cdot)$ \omterm{def:b3:lebesgue:l1}{قابلة للمكاملة} من أجل كل $t$؛
ونضع $F(t) = \int_X
f(t, x)\,\dd\mu(x)$.

\begin{theorem}[الاتصال]\label{thm:b3:lebesgue:paramcont}
لنفترض: أن $t \mapsto f(t,x)$ متصلة عند $t_0$ من أجل $x$ \omterm{def:b3:lebesgue:l1}{في كل مكان
تقريبًا}، وأن هناك \omterm{def:b3:lebesgue:l1}{دالة قابلة للمكاملة} $g$ تحقق $\abs{f(t,x)} \leq
g(x)$ من أجل
كل $t$ في جوار للنقطة $t_0$ ومن أجل $x$ \omterm{def:b3:lebesgue:l1}{في كل مكان تقريبًا}.
عندئذٍ تكون $F$ متصلة عند $t_0$.\index{التكاملات المتعلقة بوسيط}
\end{theorem}

\begin{proof}
من أجل أي متتالية $t_n \to t_0$: $f(t_n, \cdot) \to f(t_0,
\cdot)$ \omterm{def:b3:lebesgue:l1}{في كل مكان تقريبًا}،
ومهيمَن عليها بالدالة $g$: فتعطي مبرهنة التقارب المهيمن أن $F(t_n) \to F(t_0)$؛
ويكفي الاتصال بالمتتاليات في الفضاءات المترية (\cref{rem:b3:topology:sequences}).
\end{proof}

\begin{theorem}[الاشتقاق تحت التكامل]
\label{thm:b3:lebesgue:paramdiff}
لتكن $T$ فترةً مفتوحةً من $\R$. لنفترض: أن $t \mapsto f(t,x)$ قابلة
للاشتقاق على $T$ من أجل $x$ \omterm{def:b3:lebesgue:l1}{في كل مكان تقريبًا}، مع
\[
\Bigl|\frac{\partial f}{\partial t}(t, x)\Bigr| \leq g(x)
\quad \text{من أجل كل } t \in T,\ \text{في كل مكان تقريبًا من أجل } x,
\]
حيث $g$ \omterm{def:b3:lebesgue:l1}{قابلة للمكاملة}. عندئذٍ تكون $F$ قابلة للاشتقاق على $T$
مع $F'(t)
= \int_X \frac{\partial f}{\partial t}(t, x)\,\dd\mu(x)$.
\end{theorem}

\begin{proof}
لنثبّت $t$ ولنأخذ $h_n \to 0$: فتحقق نسب التزايد
\[
\varphi_n(x) = \frac{f(t + h_n, x) - f(t, x)}{h_n}
\longrightarrow \frac{\partial f}{\partial t}(t,x)
\quad\text{في كل مكان تقريبًا},
\]
وتحدّ متراجحة القيم المتوسطة المقدارَ $\abs{\varphi_n(x)} \leq
\sup_{s}\abs{\partial_tf(s,x)} \leq g(x)$: فتنطبق مبرهنة
التقارب المهيمن، و $\frac{F(t + h_n) - F(t)}{h_n} = \int\varphi_n \to
\int\partial_t f(t, \cdot)$.
\end{proof}

\begin{example}[دالة غاما]\label{ex:b3:lebesgue:gamma}
من أجل $t > 0$ نضع\index{دالة غاما}
\[
\Gamma(t) = \int_0^{+\infty} x^{t-1}\eu^{-x}\,\dd x .
\]
ويتقارب التكامل: فقرب $0$ تكون $x^{t-1}$ \omterm{def:b3:lebesgue:l1}{قابلة للمكاملة}
($t >
0$)؛ وعند اللانهاية $x^{t-1}\eu^{-x} \leq C\eu^{-x/2}$. وتعطي المكاملة بالتجزئة (على
$[\varepsilon, A]$، ثم بالنهايات عبر مبرهنة التقارب الرتيب) المعادلةَ
الدالية $\Gamma(t + 1) =
t\,\Gamma(t)$، ومنه $\Gamma(n+1) = n!$: أي العاملي مستوفًى. وعلى كل
$\intcc ab \subseteq \intoo0{+\infty}$، يكون $\partial_t\bigl(x^{t-1}\eu^{-x}\bigr) = \ln
x\cdot x^{t-1}\eu^{-x}$ مهيمَنًا عليه بالمقدار $\abs{\ln x}(x^{a-1} +
x^{b-1})\eu^{-x}$ القابل للمكاملة:
ومنه فإن $\Gamma$ من الصنف $\mathcal C^1$، وبالتراجع $\mathcal C^\infty$، مع
$\Gamma^{(k)}(t) =
\int_0^\infty(\ln x)^kx^{t-1}\eu^{-x}\dd x$. والقيمة $\Gamma(\frac12) = \sqrt\pi$ هي التكامل الغاوسي متنكّرًا
(\cref{pb:b3:lebesgue:1}).
\end{example}

\begin{omfigure}
\begin{tikzpicture}
\begin{axis}[omaxis, width=10.6cm, height=5.4cm,
    xmin=0, xmax=25, ymin=-0.25, ymax=1.05,
    xtick={0, 3.1416, 6.2832, 9.4248, 12.566, 15.708, 18.850,
      21.991, 25.133},
    xticklabels={$0$, $\pi$, $2\pi$, $3\pi$, $4\pi$, $5\pi$,
      $6\pi$, $7\pi$, $8\pi$},
    ytick={0, 1}]
  \addplot[omDef!40] coordinates {(0,0) (25,0)};
  \addplot[omThm, thick, domain=0.02:25, samples=600]
    {sin(deg(x))/x};
  \addplot[omDef, dashed, domain=0.02:25, samples=200]
    {1/x};
  \addplot[omDef, dashed, domain=0.02:25, samples=200]
    {-1/x};
\end{axis}
\end{tikzpicture}

{\small المقدار المكامَل $\frac{\sin x}x$: يتقارب التكامل
المعتلّ $\int_0^\infty$ بالتلاشي المتناوب بين الأقواس، لكن
مساحات $\abs{\cdot}$ للأقواس تسلك سلوك $\frac2{\pi k}$ --- أي
متسلسلة توافقية: $\frac{\sin x}x
\notin L^1$. فالقابلية للمكاملة بمعنى لوبيغ هي
القابلية للمكاملة المطلقة.}
\end{omfigure}

\section{تمارين}

\begin{exercise}[$\star$]\label{exo:b3:lebesgue:1}
(a) برهن على أن كل دالة رتيبة $\R \to \R$ قابلة \omterm{def:b3:measure:measure}{للقياس} بمعنى
بوريل، وعلى أن مشتق دالة قابلة للاشتقاق في كل نقطة قابلٌ \omterm{def:b3:measure:measure}{للقياس}
بمعنى بوريل.
(b) برهن على أن $f \colon X \to \R$ قابلة \omterm{def:b3:measure:measure}{للقياس} إذا وفقط إذا كان $\{f > q\}
\in \mathcal A$ من
أجل كل عدد \emph{ناطق} $q$.
\end{exercise}

\begin{exercise}[$\star$]\label{exo:b3:lebesgue:2}
احسب، بتبرير كامل:
\[
\lim_{n\to\infty}\int_0^{+\infty}
\frac{\cos x}{(1 + x/n)^{n}}\,\dd x,
\qquad
\lim_{n\to\infty}\int_0^1 \frac{n\,x^{n-1}}{1 + x}\,\dd x .
\]
\emph{(من أجل الثاني: عوّض $u = x^n$ قبل الهيمنة.)}
\end{exercise}

\begin{exercise}[$\star\star$]\label{exo:b3:lebesgue:3}
(a) أعطِ مثالًا على متراجحة تامة في مبرهنة فاتو المساعدة.
(b) أعطِ $f_n \to 0$ نقطةً نقطة مع $\int f_n = 1$ بثلاث طرائق:
بالهروب في الارتفاع، وفي العرض، وإلى اللانهاية. وأي فرضية واحدة
من فرضيات مبرهنة التقارب المهيمن تخالفها كلٌّ منها؟
(c) برهن على أنه لا يمكن في مبرهنة فاتو المساعدة استبدال
$\limsup$ بالمقدار $\liminf$ في أي من الطرفين.
\end{exercise}

\begin{exercise}[$\star\star$]\label{exo:b3:lebesgue:4}
(a) برهن على $\displaystyle\int_0^{+\infty}\frac{x}{\eu^x -
1}\,\dd x = \sum_{n\geq1}\frac1{n^2} = \frac{\pi^2}6$ \emph{(انشر $\frac1{\eu^x - 1}$ في متسلسلة هندسية
وكامل حدًّا حدًّا --- وأي مبرهنة تسمح بذلك؟)}.
(b) (حلم طالب السنة الثانية) برهن على $\displaystyle\int_0^1 x^{-x}\,\dd x = \sum_{n\geq1}n^{-n}$. \emph{(اكتب
$x^{-x} = \eu^{-x\ln x} = \sum_k\frac{(-x\ln
x)^k}{k!}$ واحسب $\int_0^1(-x\ln x)^k\dd x$ بالتعويض $x = \eu^{-u/(k+1)}$، متعرّفًا على $\Gamma$.)}
\end{exercise}

\begin{exercise}[$\star\star$]\label{exo:b3:lebesgue:5}
(a) برهن على أن $f \geq 0$ القابلة \omterm{def:b3:measure:measure}{للقياس} والمحققة $\int f\,\dd\mu = 0$ تحقق
$f = 0$ \omterm{def:b3:lebesgue:l1}{في كل مكان تقريبًا}. \emph{(انظر في $\{f \geq 1/n\}$
ومتراجحة ماركوف: $\mu(\{f \geq a\}) \leq \frac1a\int f$ --- وبرهن عليها.)}
(b) برهن على أن \omterm{def:b3:lebesgue:l1}{الدالة القابلة للمكاملة} $f$ منتهية \omterm{def:b3:lebesgue:l1}{في كل مكان
تقريبًا}.
(c) برهن على أنه إذا كان $\int_A f\,\dd\mu = 0$ من أجل \emph{كل} مجموعة قابلة
\omterm{def:b3:measure:measure}{للقياس} $A$، فإن $f = 0$ \omterm{def:b3:lebesgue:l1}{في كل مكان تقريبًا}.
\end{exercise}

\begin{exercise}[$\star\star$]\label{exo:b3:lebesgue:6}
(a) طبّق \cref{thm:b3:lebesgue:riemann} لتقرّر القابلية للمكاملة بمعنى ريمان من أجل:
$\mathbf 1_\Q$؛ ودالة توماي؛ و $\mathbf
1_K$ حيث $K$ مجموعة كانتور
سمينة (\cref{exo:b3:measure:5}).
(b) برهن على $\int_1^{+\infty}\abs{\frac{\sin x}x}\,\dd x =
+\infty$، بينما يوجد $\lim_{A\to\infty}\int_1^A\frac{\sin
x}x\,\dd x$ \emph{(كامل بالتجزئة)}:
أي تقارب معتلّ دون قابلية للمكاملة.
\end{exercise}

\begin{exercise}[$\star\star$]\label{exo:b3:lebesgue:7}
برّر أن $F(t) = \int_0^{+\infty}\eu^{-x^2}\cos(tx)\,\dd x$ من الصنف $\mathcal C^1$ على $\R$ ويحقق $F'(t) = -\frac t2F(t)$
\emph{(كامل بالتجزئة)}؛ واستنتج $F(t) =
F(0)\,\eu^{-t^2/4}$. (ومع $F(0) = \frac{\sqrt\pi}2$ من مسألة
نهاية الأسبوع: تكون الدالة الغاوسية في جوهرها تحويلَ فورييه
لنفسها --- وسينظّم \cref{ch:b3:fouriertransform} ذلك.)
\end{exercise}

\begin{exercise}[$\star\star\star$]\label{exo:b3:lebesgue:8}
(فرولاني) ليكن $0 < a < b$. برهن على
\[
\int_0^{+\infty}\frac{\eu^{-ax} - \eu^{-bx}}{x}\,\dd x =
\ln\frac ba,
\]
بكتابة المقدار المكامَل على الصورة $\int_a^b \eu^{-xt}\,\dd t$ وتبرير التبديل
بواسطة النظرية غير السالبة (\cref{cor:b3:lebesgue:additivity} في صيغتها المتصلة ---
باستباق مبرهنة تونيلي، أو بتشريح $[a,b]$ إلى $n$ جزءًا متساويًا
والانتقال إلى النهاية).
\end{exercise}

\begin{exercise}[$\star\star$]\label{exo:b3:lebesgue:9}
لتكن $f \geq 0$ قابلة \omterm{def:b3:measure:measure}{للقياس} على $(X, \mathcal A, \mu)$. برهن على أن $\nu(A) = \int_A f\,\dd\mu$
تعرّف قياسًا (\emph{كثافته}\index{كثافة قياس} $f$ بالنسبة إلى
$\mu$)، وعلى أن $\int g\,\dd\nu = \int gf\,\dd\mu$ من أجل كل $g \geq 0$ قابلة \omterm{def:b3:measure:measure}{للقياس}
\emph{(برهن عليها من أجل الدوال المميّزة، ثم البسيطة، ثم بمبرهنة
التقارب الرتيب --- وهي الآلة المعيارية)}.
\end{exercise}

\begin{exercise}[$\star\star\star$]\label{exo:b3:lebesgue:10}
(إخفاق على طريقة فايرشتراس) نعرّف $f(t) =
\int_0^{+\infty}\frac{\sin(tx)}{x(1 + x^2)}\,\dd x$.
(a) برهن على أن $f$ معرَّفة تعريفًا سليمًا ومتصلة على $\R$، ومن
الصنف $\mathcal C^1$ مع $f'(t) = \int_0^\infty\frac{\cos(tx)}{1 +
x^2}\dd x$ من أجل كل $t$ --- لكن الاشتقاق
\emph{مرة أخرى} تحت التكامل غير مشروع.
(b) بقبول $f'(t) = \frac\pi2\eu^{-t}$ من أجل $t > 0$ (المبرهَن عليه في \cref{ch:b3:residues})،
ما قيمة $\int_0^\infty\frac{x\sin(tx)}{1+x^2}\dd x$ من أجل $t > 0$، ولماذا تؤكّد صيغتها الإخفاقَ في
(a)؟
\end{exercise}

\begin{exercise}[$\star\star$]\label{exo:b3:lebesgue:11}
(مبرهنة شيفيه المساعدة) لتكن $f_n, f \geq 0$ \omterm{def:b3:lebesgue:l1}{قابلة للمكاملة} مع
$f_n \to f$ \omterm{def:b3:lebesgue:l1}{في كل مكان تقريبًا} و $\int f_n \to \int f$.
(a) برهن على $\int\abs{f_n - f} \to 0$. \emph{(طبّق التقارب المهيمن على $g_n = (f - f_n)^+ \leq f$،
واكتب $\int\abs{f_n - f} = 2\int g_n - \int(f - f_n)$.)}
(b) برهن بمثال على أنه لا يمكن إسقاط الفرضية $\int f_n \to \int
f$ (بنتوء
منزلق أو متركّز)، وعلى أن النتيجة تخفق من أجل $f_n$ ذات إشارة
دون تحكّم بالقيمة المطلقة: فإن $f_n = n\mathbf 1_{\intoc0{1/n}} -
n\mathbf 1_{\intoc{-1/n}0}$ يحقق $f_n \to 0$ \omterm{def:b3:lebesgue:l1}{في كل
مكان تقريبًا} و $\int f_n
= 0 \to 0$، ومع ذلك $\int\abs{f_n} = 2$.
(c) تطبيق (الكثافات): إذا كانت كثافات احتمالية تحقق $p_n
\to p$ \omterm{def:b3:lebesgue:l1}{في
كل مكان تقريبًا}، فإن $\int\abs{p_n - p} \to 0$ تلقائيًا: أي إن التقارب نقطةً نقطة
للكثافات هو تقارب في $L^1$ --- ترقيةٌ مجانية للتقارب.
\end{exercise}

\begin{exercise}[$\star\star$]\label{exo:b3:lebesgue:12}
نهايات كلاسيكية، بتبرير كامل عبر مبرهنتَي التقارب
الرتيب/المهيمن:
\[
\text{(a)}\ \lim_{n\to\infty}\int_0^n\Bigl(1 -
\frac xn\Bigr)^n\eu^{x/2}\,\dd x, \qquad
\text{(b)}\ \lim_{n\to\infty}\int_0^1\frac{n\,x^{n-1}}{1 +
x}\,\dd x,
\]
\[
\text{(c)}\ \lim_{n\to\infty}\int_0^\infty
\frac{\dd x}{(1 + x/n)^n\,x^{1/n}} .
\]
\emph{(من أجل (a): $(1 - x/n)^n \nearrow \eu^{-x}$ من أجل $x$ مثبَّت --- برهن على
الرتابة عبر $\log$؛ ومن أجل (b)، كامل بالتجزئة أو عوّض
$x = u^{1/n}$ وعيّن تركّزًا حدّيًا؛ ومن أجل (c)، جد مهيمِنًا
قابلًا للمكاملة صالحًا من أجل كل $n \geq 2$ بالتشريح عند
$x =
1$.)}
\end{exercise}

\section{مسألة: تكاملان شهيران}

\begin{problem}[{مسألة نهاية الأسبوع --- التكامل الغاوسي وتكامل
ديريكليه، بالوسائط وحدها}]
\label{pb:b3:lebesgue:1}
تكاملان يحكمان التحليل التطبيقي:
\[
G = \int_{-\infty}^{+\infty}\eu^{-x^2}\dd x = \sqrt\pi,
\qquad
D = \int_0^{+\infty}\frac{\sin x}{x}\,\dd x = \frac\pi2
\]
(والثاني بوصفه تكاملًا معتلًّا، \cref{ex:b3:lebesgue:riemannexamples}). ونبرهن عليهما
باستعمال أدوات هذا الفصل وحدها.

\medskip
\noindent\textbf{الجزء الأول --- التكامل الغاوسي.} من أجل
$t \geq 0$ نضع
\[
A(t) = \Bigl(\int_0^t\eu^{-x^2}\dd x\Bigr)^{2},
\qquad
B(t) = \int_0^1\frac{\eu^{-t^2(1 + x^2)}}{1 + x^2}\,\dd x .
\]
\begin{enumerate}
  \item برّر أن $A$ و $B$ من الصنف $\mathcal C^1$ على $\intoo0{+\infty}$
        واحسب $A'$ و $B'$؛ وبرهن على $A'(t) + B'(t) = 0$. \emph{(في $B'$،
        عوّض $u =
        tx$.)}
  \item احسب $A(0) + B(0)$ و $\lim_{t\to+\infty}(A + B)(t)$ --- وبرّر النهاية تحت
        التكامل في $B$.
  \item اخلص إلى $\int_0^\infty \eu^{-x^2}\dd x =
        \frac{\sqrt\pi}2$، ومنه $G = \sqrt\pi$، واستنتج $\Gamma(\tfrac12) = \sqrt\pi$
        \emph{(عوّض $x =
        u^2$ في $\Gamma(\frac12)$)}.
\end{enumerate}

\medskip
\noindent\textbf{الجزء الثاني --- تكامل ديريكليه.} من أجل
$t
\geq 0$ نضع
\[
F(t) = \int_0^{+\infty}\eu^{-tx}\,\frac{\sin x}{x}\,\dd x .
\]
\begin{enumerate}[resume]
  \item برهن على أن التكامل المعرِّف للدالة $F(t)$ يتقارب من أجل كل
        $t > 0$ بوصفه تكامل لوبيغ، ومن أجل $t = 0$ بوصفه تكاملًا
        معتلًّا؛ وبرهن على وجود $D =
        \lim_{A\to\infty}\int_0^A\frac{\sin x}x\dd x$ \emph{(كامل بالتجزئة على
        $[\pi, A]$)}.
  \item برهن على أن $F$ من الصنف $\mathcal C^1$ على $\intoo0{+\infty}$ مع
        \[
        F'(t) = -\int_0^{+\infty}\eu^{-tx}\sin x\,\dd x
        = -\frac{1}{1 + t^2}
        \]
        \emph{(بالهيمنة على $[t_0, \infty)$ من أجل كل $t_0 >
        0$؛
        والتكامل الأخير بمكاملتين بالتجزئة أو بالأُسّيات
        العقدية)}.
  \item برهن على أن $F(t) \to 0$ حين $t \to +\infty$، واستنتج
        $F(t)
        = \frac\pi2 - \arctan t$ على $\intoo0{+\infty}$.
  \item النقطة الدقيقة: $D = \lim_{t\to0^+}F(t)$. برهن عليها بتحكّم منتظم في
        الذيل: فمن أجل $0 \leq t \leq 1$ ومن أجل $A \geq \pi$، كامل
        بالتجزئة لتبرهن على
        \[
        \Bigl|\int_A^{+\infty}\eu^{-tx}\frac{\sin
        x}x\,\dd x\Bigr| \leq \frac{C}{A}
        \]
        حيث $C$ لا يتعلق بالوسيط $t$ \emph{(اشتق $\frac{\eu^{-tx}}x$ وحُدّ
        $\abs{\cos}$ بالعدد $1$؛ ولاحظ $t\eu^{-tx} \leq 1/x\cdot(tx\eu^{-tx})$ مع $\sup_{u\geq0}u\eu^{-u} < 1$)}؛ ثم
        شرّح $F(t) - D$ إلى $[0, A]$ (حيث تنطبق مبرهنة التقارب
        المهيمن حين $t \to 0$) وإلى $[A, \infty)$.
  \item اخلص إلى: $D = \frac\pi2$.
\end{enumerate}

\medskip
\noindent\textbf{الجزء الثالث --- العوائد.}
\begin{enumerate}[resume]
  \item احسب $\int_0^{+\infty}\frac{\sin^2x}{x^2}\,\dd x$ \emph{(كامل بالتجزئة واختزل إلى $D$ عبر
        $\sin
        2x = 2\sin x\cos x$)}.
  \item احسب $\int_0^{+\infty}\frac{1 -
        \cos x}{x^2}\,\dd x$، وتحقق من اتساق النتيجتين.
  \item من أجل $a > 0$، احسب $\int_0^{+\infty}\frac{\sin(ax)}x\dd x$ و $\int_{-\infty}^{+\infty}\eu^{-ax^2}\dd x$، وسجّل قواعد
        التحجيم (فستكون هي أحصنة العمل في \cref{ch:b3:fouriertransform}).
  \item فسّر بدقة لماذا لم يكن ممكنًا معالجة $D$ بمبرهنة التقارب
        المهيمن مباشرةً عند $t = 0$ (فلا مهيمِن قابل للمكاملة
        على $[0,1]\times[0,\infty)$)، ولماذا يكون تشريح الذيل في السؤال 7 هو
        البديل الأمين --- فهذا النمط («القابلية المنتظمة
        لمكاملة الذيول») يتكرّر في التحليل كله.
\end{enumerate}

\medskip
\noindent\textbf{الجزء الرابع --- \omterm{ex:b3:lebesgue:gamma}{دالة غاما} حسب بور
ومولروب.} تحقق الدالة $\Gamma$ (\cref{ex:b3:lebesgue:gamma}) الشرطَ
$\Gamma(1) = 1$ و $\Gamma(x+1) = x\Gamma(x)$ --- لكن عددًا لامنتهيًا من الدوال
الأخرى تحققه أيضًا (بالضرب في أي تموّج دوري بالدور $1$). ويحدّد
شرط تحدّب واحد دالةَ $\Gamma$ تحديدًا وحيدًا، ثم تسقط متطابقاتها
الأعمق آليًا. ونقول عن دالة موجبة $f$ على فترة إنها
\emph{محدّبة لوغاريتميًا} إذا كانت $\log f$ محدّبة.
\begin{enumerate}[resume]
  \item برهن على أن التحدّب اللوغاريتمي يستلزم التحدّب، وعلى أن
        جداءات الدوال المحدّبة لوغاريتميًا ومركّباتها مع
        التطبيقات التآلفية محدّبة لوغاريتميًا، وعلى أن --- بواسطة
        متراجحة هولدر ذات الدالتين $\int\abs{uv} \leq
        \bigl(\int\abs u^p\bigr)^{1/p}\bigl(\int\abs
        v^q\bigr)^{1/q}$، المبرهَن عليها
        مباشرةً من متراجحة يونغ --- الدالةَ $\Gamma$ محدّبة
        لوغاريتميًا على $\intoo0\infty$.
  \item (مبرهنة الميل المساعدة) لتكن $g$ محدّبة على
        $\intoo0\infty$ مع $g(n+1) - g(n) = \log n$ من أجل كل عدد صحيح $n
        \geq 1$.
        من أجل $x \in \intoc01$ و $n \geq 2$، قارن ميول $g$ على
        $[n-1, n]$ و $[n, n+x]$ و $[n, n+1]$، واستنتج
        \[
        x\log(n-1) \;\leq\; g(n + x) - g(n) \;\leq\; x\log n .
        \]
  \item (بور--مولروب) لتكن $f > 0$ تحقق $f(1) = 1$
        و $f(x+1) = xf(x)$ و $\log f$ محدّبة. بفكّ العلاقة
        التراجعية إلى $f(n + x) = x(x+1)\cdots(x + n -
        1)\,f(x)$ و $f(n) = (n-1)!$، استنتج من
        السؤال 14 أنه من أجل $x \in \intoc01$
        \[
        f(x) = \lim_{n\to\infty}
        \frac{n!\,n^x}{x(x+1)\cdots(x+n)} :
        \]
        فتكون $f$ وحيدة، ومنه $f = \Gamma$، وتصح \emph{صيغة
        غاوس للنهاية} (وتُمدَّد إلى كل $x > 0$ بالعلاقة
        التراجعية).
  \item نعرّف \emph{دالة بيتا} $B(x, y) =
        \int_0^1t^{x-1}(1-t)^{y-1}\,\dd t$ (حيث $x, y > 0$). برهن
        على التقارب، وعلى العلاقة التراجعية $B(x+1, y) =
        \frac{x}{x+y}\,B(x, y)$ (بالمكاملة
        بالتجزئة)، وعلى $B(1, y) = \frac1y$.
  \item برهن على أن $x \mapsto B(x, y)$ محدّبة لوغاريتميًا (بهولدر مرة
        أخرى)، وطبّق بور--مولروب على
        \[
        f(x) = \frac{B(x, y)\,\Gamma(x + y)}{\Gamma(y)}
        \]
        لتخلص إلى \emph{صيغة أويلر}: $B(x, y) =
        \dfrac{\Gamma(x)\Gamma(y)}{\Gamma(x+y)}$ --- دون أي تكاملات
        مضاعفة.
  \item احسب $B(\frac12, \frac12)$ مباشرةً (عوّض $t
        = \sin^2\theta$) واستنتج $\Gamma(\frac12) =
        \sqrt\pi$: أي
        التكامل الغاوسي من الجزء الأول، مستعادًا بمحض التحدّب.
        وقارن البرهانين بجملة واحدة لكلٍّ منهما.
  \item (تضعيف لوجندر) برهن على أن
        \[
        g(x) = \frac{2^{x-1}}{\sqrt\pi}\,
        \Gamma\Bigl(\frac x2\Bigr)
        \Gamma\Bigl(\frac{x+1}2\Bigr)
        \]
        تحقق فرضيات بور--مولروب الثلاث، واخلص إلى $g = \Gamma$،
        أي $\Gamma(2z) =
        \frac{2^{2z-1}}{\sqrt\pi}\,\Gamma(z)\,\Gamma(z +
        \tfrac12)$ من أجل كل $z > 0$.
  \item استنتج الصيغة المغلقة $\Gamma\bigl(n + \tfrac12\bigr)
        = \dfrac{(2n)!}{4^n\,n!}\sqrt\pi$، وبرهن، بتطبيق مبرهنة
        الميل المساعدة على $\log\Gamma$ حول الأعداد الصحيحة
        الكبيرة، على المقاربات
        \[
        \frac{\Gamma(n + \frac12)}{\Gamma(n)\,\sqrt n}
        \longrightarrow 1 .
        \]
  \item اجمع السؤالين الأخيرين في مقاربات المعامل الثنائي
        المركزي
        \[
        \binom{2n}{n} \sim \frac{4^n}{\sqrt{\pi n}},
        \]
        وتحقق عدديًا من أجل $n = 10$ ($\binom{20}{10} =
        184756$، مقابل $4^{10}/\sqrt{10\pi} \approx
        187079$:
        فالنسبة $\approx 0.988$).
  \item (تركيب) ظهر الثابت $\sqrt\pi$ حتى الآن بوصفه التكامل
        الغاوسي (الجزء الأول)، وبوصفه $B(\frac12,
        \frac12)$ (السؤال 18)،
        وداخل التضعيف (السؤال 19)؛ ويستبق تقدير المعامل الثنائي
        المركزي كلًّا من ستيرلنغ (مسألة نهاية الأسبوع في
        \cref{ch:b3:product}) ودي موافر--لابلاس. ارسم الصلات: أي العبارات
        تكافئ أيها، وماذا تسهم به كل تقنية --- الاشتقاق تحت
        التكامل في مقابل التحدّب --- مما لا تستطيعه الأخرى؟
\end{enumerate}

\medskip
\noindent\textbf{الجزء الخامس --- ثلاثة عوائد أخرى.}
\begin{enumerate}[resume]
  \item (واليس، بدالة بيتا) من أجل $p > -1$، عوّض $t =
        \sin^2\theta$ لتبرهن
        على
        \[
        W_p = \int_0^{\pi/2}\sin^p\theta\,\dd\theta
        = \frac12\,B\Bigl(\frac{p+1}2, \frac12\Bigr),
        \]
        واستنتج من العلاقة التراجعية لدالة بيتا (السؤال 16) أن
        $W_{n+2} = \frac{n+1}{n+2}\,W_n$ من أجل الأعداد الصحيحة $n \geq
        0$. واحسب $W_{2n}$
        و $W_{2n+1}$ بصيغة مغلقة، وبرهن على $W_{2n+1}/W_{2n} \to 1$ بالحصر،
        واخلص إلى \emph{جداء واليس}
        \[
        \frac\pi2 = \lim_{n\to\infty}
        \prod_{k=1}^{n}\frac{4k^2}{4k^2 - 1} .
        \]
  \item (الدالة الغاوسية تلتقي تواترًا) من أجل $b \in \R$ نضع
        \[
        \Phi(b) = \int_{-\infty}^{+\infty}
        \eu^{-x^2}\cos(2bx)\,\dd x .
        \]
        برهن على أن $\Phi$ من الصنف $\mathcal C^1$ على $\R$، وعلى
        أن مكاملة بالتجزئة تعطي المعادلة التفاضلية $\Phi'(b) = -2b\,\Phi(b)$،
        واخلص إلى
        \[
        \Phi(b) = \sqrt\pi\,\eu^{-b^2} :
        \]
        أي إن الدالة الغاوسية تُعيد إنتاج نفسها بهذا التحويل ---
        وهي المتطابقة الوحيدة التي سيدور عليها \cref{ch:b3:fouriertransform}.
  \item (تكامل فرولاني) من أجل $0 < a < b$، برهن على
        \[
        \int_0^{+\infty}
        \frac{\eu^{-ax} - \eu^{-bx}}{x}\,\dd x
        = \log\frac ba ,
        \]
        بالاشتقاق بالنسبة إلى الوسيط $a$ (وبرّر الهيمنة على كل
        $\intco{a_0}{+\infty}$ حيث $a_0 > 0$، وعيّن الثابت بجعل $a \to b$).
        وأين يحتاج المقدار المكامَل بالضبط إلى تفرّده القابل
        للإزالة عند $x = 0$؟
\end{enumerate}
\end{problem}
